\noindent
Large language models (LLMs) are flooding into Human-Centered Computing (HCC) research as substitutes for human
judgment: coding qualitative data, annotating crowdsourced corpora, simulating study participants, generating
personas, and evaluating designs. Yet HCC's methods carry epistemological commitments (interpretation, reflexivity,
positionality) that make uncritical substitution risky, and the field lacks an evidence-based account of
\emph{when} an ``LLM-as-a-judge'' can stand in for a human and when it cannot. We address this with four studies.
\textbf{Study 1a} is a PRISMA systematic review of \Nstudies{} studies across \Ndomains{} domains that establishes
the general reliability landscape of LLM-as-a-judge: substitution is validated ($\approx$ or above the
human-human ceiling) for objective, verifiable, low-expertise tasks judged in aggregate, and unreliable for
subjective, expert, low-resource, or adversarial ones (mean reliability \Meanrel/5; only \Pctreliable\% of studies
in the reliable band). \textbf{Study 1b} asks whether this landscape is a moving target: we re-run a flagship
``promising'' benchmark (SummEval) with the most powerful open models served via Clemson's vLLM API and ask whether
they now \emph{surpass} the agreement numbers reported for proprietary judges, distinguishing a ``matter of time''
from a structural ceiling. \textbf{Study 2} is a PRISMA systematic review of \STN{} studies focused on HCC, HCI, and
qualitative-research methods that historically require human input. It finds that the field overwhelmingly positions
LLMs to \emph{augment} rather than replace humans (\STaugment{} augment vs.\ \STreplace{} replace vs.\ \STsimulate{}
simulate), with an epistemic stance dominated by \emph{tension} (\STtension/\STN); objective crowd-annotation
substitution is validated, interpretive qualitative analysis is reliable only as human-in-the-loop augmentation, and
\emph{simulating} participants, personas, and usability evaluations is the least reliable (flattening of identity
groups, caricature, hallucinated usability issues). \textbf{Study 3} closes the gaps the
first two studies expose, by running a standardized LLM-as-a-judge meta-evaluation, using open models served via
Clemson's vLLM API and the agreement-against-human-ceiling methodology distilled from Studies 1-2, on HCC datasets
that contain human labels but have never been formally tested with an LLM judge. Across \STHRn{} judgments (a
6-model open panel, 9B-754B), it finds the human-agreement ceiling governs reliability: the panel matches or
exceeds humans on low-ceiling NLI but falls short of, and fails the alt-test for, subjective social-computing
annotation (emotion, hate speech), with binary offensiveness the strongest case. A mixed-effects model confirms the
mechanism: judge correctness is governed overwhelmingly by item-level human (dis)agreement (random-intercept
$\mathrm{SD}_{\text{item}}=\SMEvcitem$) and by task subjectivity ($\mathrm{OR}=\SMEsubjOR$ per unit of
$1\!-\!$ceiling), while model scale contributes only modestly ($\mathrm{OR}=\SMEsizeOR$ per $\log_{10}$ params);
after FDR correction \SMEfdrbelow{} of \SMEfdrtests{} judge$\times$dataset cells sit significantly below the human
ceiling. Together the studies yield a
decision framework telling HCC researchers when an LLM judge is a validated substitute, a human-in-the-loop
assistant, or an unsafe replacement.
